\begin{abstract}
Individual event probabilities on temporal operational graphs do not determine
the uncertainty of group counts or simultaneous delays. We build a
current-cohort forecasting benchmark on a compute-cluster trace and two
loan-workflow releases, separating marginal prediction, shared dependence,
typed incidence and group-risk evaluation. Controls include temporal neural
and tree predictors, classical factor scenarios, direct count distributions,
historical residuals and discrete probabilistic reconciliation. Original
frozen evaluations motivate explicitly retrospective calibration and graph
diagnostics; a separately frozen workflow period provides fresh confirmation.
Typed factors reduce cluster count CRPS by 24.3\% against independence, but
only 2.2\% beyond a global factor. Recent same-version calibration and direct
count fitting substantially change the workflow comparison. In fresh
confirmation, global dependence improves count CRPS by 12.7\% while worsening
a reserve-loss proxy by 20.4\%; a development-selected direct count model
outperforms it on both. Count and peak-group-delay rankings also differ.
The contribution is a reproducible application benchmark and controlled
empirical evidence for evaluating mean bias, collective uncertainty and
operational risk separately, with explicit limits on temporal replication
and deployment interpretation.
\end{abstract}
